\section{Diagnostics and Compiler Invariants}\label{sec:compiler-diagnostics}

Having described parsing, semantic analysis, lowering, and MIDI serialization, this section ties the pipeline together:
how errors are collected across stages and the hard limits the implementation enforces to keep compilation predictable.

\subsection{Shared Diagnostics Object and Stage Gating}\label{detail-subsec:diagnostics-architecture}

Diagnostics are centralized in the \texttt{motivo\_diagnostics} library.
\texttt{motivo::compile} constructs one \texttt{DiagnosticsEngine} on the stack and passes a non-owning reference into
\texttt{parsing::parse\_stream}, \texttt{semantic::analyze}, \texttt{lowering::lower}, and the \texttt{write\_midi}
exception handler.
The engine appends \texttt{Diagnostic} records to an internal vector; no stage clones or replaces the object, so a
single compilation accumulates all messages in emission order.

Each \texttt{Diagnostic} carries a \texttt{DiagnosticStage} (\texttt{Parsing}, \texttt{Semantic}, \texttt{Lowering},
\texttt{Output}), a \texttt{DiagnosticSeverity}, an optional stringified \texttt{source::Location}, and a free-form
message.
The public formatter \texttt{format\_diagnostic} renders compiler-native output as
\texttt{error: \textless stage\textgreater: \textless location\textgreater: \textless message\textgreater}, where
stage strings are \texttt{parse}, \texttt{semantic}, \texttt{lowering}, and \texttt{output}.
Motivo Studio's compile API serializes the same records into JSON with a \texttt{type} field mirroring the stage enum;
the thesis uses the string form below because it matches unit-test assertions and CLI stderr.

\topic{Stage gating in the compiler driver.}
Compilation is strictly sequential; later stages never run on input that failed an earlier gate:

\begin{enumerate}
  \item \textbf{Parsing.} If \texttt{ParseResult::ok()} is false, \texttt{compile} returns immediately with accumulated
    parse diagnostics. Semantic analysis does not run on a null program.
  \item \textbf{Semantic.} If \texttt{diagnostics.has\_errors(DiagnosticStage::Semantic)}, lowering is skipped even when
    the AST is structurally complete.
  \item \textbf{Lowering.} If \texttt{LowerResult::ok()} is false, MIDI serialization is skipped.
  \item \textbf{Output.} \texttt{write\_midi} exceptions are caught and reported as \texttt{DiagnosticStage::Output}
    errors; prior stage diagnostics are retained.
\end{enumerate}

\topic{Lexical error collection.}
The Flex scanner calls \texttt{report\_lexical}, which forwards to
\texttt{DiagnosticsEngine::report} when a diagnostics pointer is installed.
Lexical errors do not abort tokenization immediately; the scanner continues so multiple invalid characters or
unterminated comments can be reported in one pass.

\topic{Parser error recovery.}
Bison is configured with \texttt{\%define parse.error detailed} and \texttt{\%define parse.lac full}.
Statement-level productions include an \texttt{error} alternative that synchronizes on the next statement terminator
(\texttt{;}), discards the erroneous fragment, and calls
\texttt{yyerrok}, and resumes at the next statement boundary.
This allows the parser to collect more than one syntax diagnostic per file while still returning a null
\texttt{ParseResult} when the overall parse fails.

\topic{Semantic and lowering collection.}
Semantic errors flow through \texttt{Traversal::diagnose}; the traversal does not short-circuit on the first error, so
independent violations (e.g., two incompatible assignments in separate tracks) surface together.
Lowering catches \texttt{LoweringFailure} exceptions, converts them to
\texttt{DiagnosticStage::Lowering} entries, and uses \texttt{register\_events} to enforce the global event cap before
MIDI generation.

\subsection{Diagnostic Taxonomy and Sample Messages}\label{detail-subsec:diagnostic-taxonomy}

Table~\ref{tab:diagnostic-taxonomy} groups the most common diagnostic families.
Listing~\ref{lst:diagnostic-samples} shows verbatim messages exercised by the Google Test suite.

\begin{table}[htbp]
  \centering
  \renewcommand{\arraystretch}{1.25}
  \begin{tabularx}{\textwidth}{@{}M{1.8cm} M{3.6cm} >{\raggedright\arraybackslash\hspace{0pt}}X@{}}
    \toprule
    \textbf{Stage} & \textbf{Origin} & \textbf{Typical trigger} \\
    \midrule
    Parsing & Flex \texttt{report\_lexical} & Invalid character, unterminated block comment \\
    Parsing & Bison \texttt{parse.error detailed} & Missing semicolon, malformed header, duplicate \texttt{tempo} \\
    Semantic & \texttt{validate\_header} & Tempo or time signature outside allowed range \\
    Semantic & \texttt{validate\_call} & Undefined pattern, overload mismatch, incompatible argument type \\
    Semantic & \texttt{validate\_binary\_operands} & Non-numeric arithmetic, non-boolean logic, invalid \texttt{==} \\
    Semantic & \texttt{visit\_*} visitors & Undeclared variable, illegal outer-scope assignment, drum/melody mismatch \\
    Lowering & \texttt{LoweringFailure} & Division or modulo by zero at runtime \\
    Lowering & \texttt{register\_events} & Total emitted \texttt{NoteEvent}s exceed \texttt{MAX\_EVENTS} \\
    Output & \texttt{write\_midi} catch & File open failure, I/O error during binary write \\
    \bottomrule
  \end{tabularx}
  \caption{Diagnostic taxonomy by pipeline stage and reporting site.}\label{tab:diagnostic-taxonomy}
\end{table}

\begin{lstlisting}[caption={Representative compiler diagnostic messages (format and wording from unit tests).},
label={lst:diagnostic-samples}]
error: semantic: <test>:1:7: tempo must be between 1 and 1000, got 1001
error: semantic: <test>:2:14: undefined variable 'x'
error: semantic: <test>:5:12: no matching overload of pattern 'motif'
error: lowering: 2:12-15: division by zero
error: lowering: 8:5-17: program exceeds 20000 event limit (has 20001)
\end{lstlisting}

\subsection{Compiler Invariants}\label{detail-subsec:compiler-invariants}

Beyond type rules, the implementation enforces several numeric and resource invariants.
Semantic checks catch invalid metadata before any note is emitted; lowering checks catch runtime failures and bound
unrolling cost.

\topic{Tempo range (semantic).}
\texttt{Traversal::validate\_header} requires \texttt{tempo} declarations to satisfy $1 \leq \mathrm{BPM} \leq 1000$.
The grammar already restricts tempo operands to positive integer literals, so only zero and values above the cap reach
semantic validation.
The upper bound prevents absurd MIDI tempo meta-events ($uspb = 60{,}000{,}000 / \mathrm{BPM}$ would underflow or
quantize poorly for extreme values) and is covered by unit tests on header validation.

\topic{Time signature rules (semantic).}
When a \texttt{signature} declaration is present, the numerator must satisfy $1 \leq n \leq 32$ and the denominator
must be a power of two in the same range.
Non-power-of-two denominators (e.g., \texttt{4/3}) are rejected because MIDI time-signature meta-events encode the
denominator as $\log_2 d$ in \texttt{build\_tempo\_track}; values that are not powers of two cannot be represented
faithfully.

\topic{Division and modulo by zero (lowering).}
Integer and floating-point division in the expression evaluator throw \texttt{LoweringFailure} when
the divisor evaluates to zero.
Modulo applies the same guard for integer operands.
These cases are not semantic errors because divisors may be variables whose value is unknown until unrolling; the test
suite confirms both literal (\texttt{1 / 0}) and variable (\texttt{4 / d}) paths.

\topic{Event cap \texttt{MAX\_EVENTS = 20000} (lowering).}
The cap provides a termination guarantee for programs with no statically-bounded iteration count.
\texttt{LowererContext::register\_events} increments a program-wide counter each time notes are emitted; exceeding
\texttt{MAX\_EVENTS} aborts lowering with a diagnostic naming the limit and the actual count.
Because Motivo accepts \texttt{for (;;)} and defers trip-count analysis to unrolling time, the cap is the stopping rule
for pathological or non-terminating loops.
Exactly 20{,}000 events across one or more tracks is valid; 20{,}001 triggers an error.

\topic{What is not statically bounded.}
The semantic phase does not limit loop iteration counts or pattern recursion depth.
A finite loop within the event cap succeeds; an infinite \texttt{for (;;)} terminates only when the event cap is hit.
This trade-off keeps the semantic pass single-pass and defers cost analysis to lowering, where emitted events are the
measurable unit of work.

With the pipeline, its data artifacts, and its diagnostic guarantees now accounted for, the compiler stands as a
complete technical artifact: source text in, validated Standard MIDI File out.
A compiler alone, however, is not something a composer can use.
The next chapter wraps \texttt{motivoc} in Motivo Studio, the browser-based product and operations layer that turns the
toolchain into a deployable service.
